% Chapter 7 — includable by main.tex
\chapter{Kontekstni jezici}
\label{chap:7}

\begin{quote}
\itshape
Bezkontekstni jezici ne mogu prebrojati tri stvari odjednom.
Kontekstni jezici rješavaju i taj problem --- dodavanjem trake
ograničene na duljinu ulaza. Klasa Kon je u mnogo čemu lijepa:
zatvorena je na sve skupovne operacije, a prepoznavanje je odlučivo.
\end{quote}


% ================================================================
\section{Motivacija: što ne može BK?}
% ================================================================

U poglavlju 6 dokazali smo da $L_\equiv = \{a^n b^n c^n \mid n \geq 1\}$
nije bezkontekstan. Intuitivno: potisni automat može uskladiti
broj $a$-ova i $b$-ova (stavi $a$-ove na stog, skidaj za svaki $b$),
ali jednom kad skine sve s hrpe, više ne zna koliko je $b$-ova prebrojio ---
to je točno potrebno za provjeru $c$-ova.

Rješenje: umjesto stoga koji nestaje, trebamo traku po kojoj se
možemo kretati u oba smjera. Ključno ograničenje: traka je dugačka
\emph{točno koliko i ulazna riječ}. Nema više memorije od ulaza.

% ================================================================
\section{Ograniceni automat (OA)}
% ================================================================

\begin{definicija}[Ograniceni automat]
\textbf{Ograniceni automat} (OA) je uredena sestorka
\[
  O = (Q,\, \Sig,\, \Gamma,\, \Delta,\, q_0,\, q_X),
\]
gdje su $Q$, $\Sig \subseteq \Gamma$, $q_0$, $q_X$ kao obicno, a
\[
  \Delta \subseteq Q \times \Gamma \times Q \times \Gamma \times \{-1, 1\}
\]
je relacija prijelaza. Prijelaz $(q, \alpha, p, \beta, d)$ znaci:
``u stanju $q$, citajuci $\alpha$ na poziciji $m$, predi u stanje $p$,
zapisi $\beta$ na poziciju $m$, te se pomakni na poziciju $m + d$
(ali ne izlazi s trake: $0 \leq m + d < |w|$).''

OA prihvaća riječ $w$ ako iz početne konfiguracije $(q_0, 0, w)$
može doći do neke završne konfiguracije $(q_X, m, t)$.
\end{definicija}

\begin{napomena}
Kljucne razlike od PA:
\begin{itemize}
    \item Traka je ograničena na $|w|$ znakova --- nema više.
    \item Može se kretati lijevo i desno (PA samo naprijed).
    \item Može pisati po traci (modifikacija ulaza).
    \item Kao i PA: OA je \textbf{nedeterminističan}, a
      deterministični i nedeterministični OA \emph{nisu} ekvivalentni
      (za razliku od KA) --- ovo je otvoreni problem!
\end{itemize}
\end{napomena}

\begin{primjer}[OA za $L_\equiv$]
Pseudokod OA koji prepoznaje $L_\equiv = \{a^n b^n c^n\}$:
\begin{enumerate}
    \item Označi prvi $a$ i pročitaj sva $a$ dalje. Ako nema $a$, odbij.
    \item Označi prvi $b$ iza svih $a$-ova. Ako nema $b$, odbij.
    \item Označi prvi $c$ iza svih $b$-ova. Ako nema $c$, odbij.
    \item Vrati se lijevo do zadnjeg označenog $a$.
    \item Ponovi od koraka 1 s prvim neoznačenim $a$.
    \item Kad su sva $a$ označena: provjeri da su sva $b$ i $c$
      također označena (ne smije biti neoznačenih). Prihvati.
\end{enumerate}
Svaki prolaz označi točno jedno $a$, jedno $b$ i jedno $c$,
sto osigurava $|a| = |b| = |c|$.
\end{primjer}

\subsection{Odlučivost $A_\text{OA}$}

Za razliku od Turingovih strojeva, za OA uvijek postoji konačno
mnogo konfiguracija (duljina je fiksna $|w|$), pa možemo problem
prihvaćanja riješiti obilazak grafa:

Za OA $O$ i rijec $w$, konstruiramo konacni graf ciji su vrhovi
sve moguce konfiguracije duljine $|w|$:
\[
  \text{broj vrhova} = |Q| \cdot |w| \cdot |\Gamma|^{|w|}.
\]
Problem prihvaćanja svodi se na dostupnost završnih konfiguracija
iz početne --- što rješavamo DFS-om ili BFS-om.

\begin{napomena}
Graf je ogroman (eksponencijalan u $|w|$), pa algoritam nije
polinoman. Ali je konačan, pa je problem \emph{odlučiv}.
\end{napomena}

% ================================================================
\section{Poboljsanja OA}
% ================================================================

Osnovna definicija OA je praktično nezgrapna. Pokazujemo niz
ekvivalentnih poboljšanja:

\begin{center}
\renewcommand{\arraystretch}{1.5}
\begin{tabular}{ll}
\hline
\textbf{Poboljsanje} & \textbf{Ideja simulacije} \\
\hline
Nemijenjanje trake & svaki $(q, *, p, *, d)$ postaje $|\Gamma|$ prijelaza \\
Više tragova & Kartezijev produkt abeceda tragova \\
Detekcija rubova & trag s oznakama $\{\wedge, \sqcup\}$ \\
Stajanje na mjestu ($d = 0$) & lijevo pa desno, uz provjeru ruba \\
Više završnih stanja & novo stanje $q_X$, $\eps$-prijelazi iz svakog $q \in F$ \\
Konačne varijable & produkt $Q \times S$ za varijablu iz $S$ \\
Više trakova (s gl.\ nezavisnim) & $2k$ tragova, svaka glava kao oznaka \\
\hline
\end{tabular}
\end{center}

Svako od ovih poboljšanja ekvivalentno je originalnom OA u smislu
da prepoznaje istu klasu jezika. Dokazi idu metodom \textbf{simulacije}:
poboljšani stroj uvijek može sve što može obični OA (trivijalni smjer),
a obični OA simulira poboljšanog imitirajući svaki elementarni korak.

% ================================================================
\section{Kontekstne i monotone gramatike}
% ================================================================

\begin{definicija}[Kontekstna gramatika]
Pravilo $uAv \to ubv$ je \textbf{kontekstno} ako su
$u, v \in Y^*$, $A \in V$ i $b \in Y^+$. Dakle: varijabla $A$ moze
se zamijeniti s $b$ samo u ``kontekstu'' $u \cdot v$ oko nje.

\textbf{Kontekstna gramatika} (KG) je gramatika u kojoj su sva pravila
kontekstna. Jezik je \textbf{kontekstan} ako ga generira neka KG.
\end{definicija}

\begin{definicija}[Monotona gramatika]
Pravilo $u \to v$ je \textbf{monotono} ako $|u| \leq |v|$ --- desna strana
nije kraca od lijeve. Gramatika je \textbf{monotona} (MG) ako su
joj sva pravila monotona.
\end{definicija}

\begin{propozicija}
Klase $\mathrm{KG}$, $\mathrm{MG}$ i $\mathrm{OA}$ generiraju/prepoznaju
istu klasu jezika: \textbf{kontekstne jezike} Kon.
Vrijedi $\mathrm{BK}_+ \subsetneq \mathrm{Kon}$ (svi pozitivni BK-jezici
su kontekstni, ali $L_\equiv$ je kontekstan a nije BK).
\end{propozicija}

\subsection{Monotona gramatika za $L_\equiv$}

Primjer monotone gramatike koja generira $L_\equiv$:
\[
  S \to abc \mid aSBc, \quad cB \to Bc, \quad bB \to bb.
\]
Izvod $a^n b^n c^n$ koristi pravilo $S \to aSBc$ tocno $n-1$ puta,
zatim $S \to abc$, pa ``proburbla'' svaki $B$ ulijevo pokraj $c$-ova
(pravilo $cB \to Bc$), i na kraju pretvori svaki $B$ u $b$ (pravilo $bB \to bb$).

Da su pravila monotona lako se provjeriti: sve desne strane su iste
ili dulje od lijevih. Da gramatika generira točno $L_\equiv$ dokazujemo
\textbf{invarijantom}: riječ $w$ u izvodu je ``balansirana'' ako ima jednak
broj $a$-ova, $b$-ova, $B$-ova i $c$-ova zajedno gledano.

\subsection{Pretvorba MG $\to$ KG}

Pretvorba ide u dvije faze analogno pretvorbi BKG u ChNF:
\begin{enumerate}
  \item \textbf{TERM}: znakove koji se pojavljuju uz ostalo zamijenimo
        novim varijablama (i na lijevoj i na desnoj strani!).
  \item \textbf{BIN}: pravila s $n \geq 2$ simbola s lijeve strane
        $A_1 \cdots A_n \to B_1 \cdots B_m$ (gdje $m \geq n \geq 2$)
        zamijenimo s $2n$ kontekstnih pravila koja ``postepeno''
        zamjenjuju $A_i$ s $C_i$, pa $C_i$ s $B_i$.
\end{enumerate}

% ================================================================
\section{Ekvivalentnost OA i MG}
% ================================================================

\subsection{MG $\to$ OA}

Za MG $G$, OA nedeterministički ``pogađa'' izvod riječi:
koristi drugi trag za međuvriječi, primjenjuje pravila u nedeterminis-
tičnom redoslijedu, i prihvaća kad se drugi trag poklopi s ulazom.
Ključno: jer su pravila monotona, međuvriječ nikad nije dulji od
ulaza --- dakle stane na traci.

\subsection{OA $\to$ MG}

Ovaj smjer je tehnički složeniji. Koristimo \textbf{jednostavni OA} (JOA)
koji prihvaća $w$ samo u konfiguraciji $(q_X, 0, w)$ --- jedinstvenoj
završnoj konfiguraciji.

Ideja: kodiramo konfiguracije OA kao rijeci nad $Y = Q \cup \Gamma$,
jer svaka konfiguracija $(q, m, t_0 \cdots t_{n-1})$ odgovara rijeci
$t_0 \cdots t_{m-1} \, q \, t_m \cdots t_{n-1}$.

Pravila gramatike simuliraju prijelaze OA:
\begin{itemize}
  \item Za prijelaz $(q, \alpha, p, \beta, 1)$: pravilo $q\alpha \to \beta p$.
  \item Za prijelaz $(q, \alpha, p, \beta, -1)$: pravilo $\gamma q\alpha \to p\gamma\beta$
        za svaki $\gamma \in \Gamma$.
\end{itemize}

Dodatni trik: objedinjavanje znakova (``sljepljivanje'' pomocnih
simbola s pravcim znakovima) osigurava da gramatika ostane monotona
dok uklanjamo rubne oznake.

% ================================================================
\section{Zatvorenost klase Kon}
% ================================================================

\begin{propozicija}
Klasa $\mathrm{Kon}$ zatvorena je na sve skupovne i jezicne operacije:
$\cup$, $\cap$, $^c$ (komplement), $\setminus$, $\cdot$, $^+$.
\end{propozicija}

Zatvorenost na uniju, konkatenaciju i plus slijedi identicki kao za BK
(dodavanjem novog početnog simbola).

Zatvorenost na \textbf{presjek} slijedi iz ``serijskog spajanja'' JOA:
završno stanje prvog postaje početno stanje drugog.
Budući da JOA čuva traku i poziciju, drugi automat počinje točno
u stanju u kojemu je prvi zavrsio.

Zatvorenost na \textbf{komplement} je tezi:

\begin{propozicija}[Immerman--Szelepcs\'enyi, 1988]
Komplement kontekstnog jezika je kontekstan.
\end{propozicija}

Ovo je netrivi jalan teorem (dokazan tek 1988., neovisno od dva
tako sto pokaze da \emph{ne postoji} prihvacajuca grana za $w$.
autora). Ideja: nedeterministični OA koji prepoznaje komplement
koristi tehniku ``brojenja'' --- nedeterministički prihvaća $w \notin L$
tako što pokaže da \emph{ne postoji} prihvaćajuća grana za $w$.
tako sto pokaze da \emph{ne postoji} prihvacajuca grana za $w$.

% ================================================================
\section{Usporedba klasa}
% ================================================================

\begin{center}
\renewcommand{\arraystretch}{1.6}
\begin{tabular}{lllll}
\hline
& \textbf{Reg} & \textbf{BK} & \textbf{Kon} & \textbf{RE} \\
\hline
Automat & KA & PA & OA & TP \\
Gramatika & DLG/LLG & BKG & KG/MG & OG \\
  Memorija & ništa & stog & ogr.\ traka & neogr.\ traka \\
Zatvorena $\cup$ & \checkmark & \checkmark & \checkmark & \checkmark \\
Zatvorena $\cap$ & \checkmark & \texttimes & \checkmark & \checkmark \\
Zatvorena $^c$ & \checkmark & \texttimes & \checkmark & \texttimes \\
 \(A_L\) odlučiv? & \checkmark & \checkmark & \checkmark & samo poluodl. \\
 \(E_L\) odlučiv? & \checkmark & \checkmark & \texttimes & \texttimes \\
 \(\mathrm{EQ}_L\) odlučiv? & \checkmark & \texttimes & \texttimes & \texttimes \\
\hline
\end{tabular}
\end{center}

\begin{zadatak}
\begin{enumerate}[label=(\alph*), itemsep=4pt]
  \item Konstruirajte OA koji prepoznaje $\{ww^R \mid w \in \{a,b\}^*\}$.
  \item Napišite monotonu gramatiku za $\{a^n b^n c^n d^n \mid n \geq 1\}$.
  \item Dokazite da je $\{a^{2^n} \mid n \geq 0\}$ kontekstan jezik.
  \item Zasto klasa Kon nije zatvorena na Kleenevu zvijezdu?
        Koji je problem s $\eps$?
\end{enumerate}
\end{zadatak}

% ================================================================
\section*{Sazetak}
% ================================================================

\medskip\noindent
\begin{tabular}{ll}
  OA & ograniceni automat: traka duljine $|w|$, kretanje lijevo-desno \\
  JOA & jednostavni OA: jedinstven zavrsni uvjet \\
  KG & kontekstna gramatika: $uAv \to ubv$ \\
  MG & monotona gramatika: $|u| \leq |v|$ \\
  $\mathrm{OA} \equiv \mathrm{KG} \equiv \mathrm{MG}$ & svi opisuju klasu Kon \\
  I-S teorem & $\mathrm{Kon}$ zatvorena na komplement \\
\end{tabular}

\medskip\noindent
U sljedećem poglavlju prelazimo na \textbf{Turingove strojeve} ---
strojeve bez ikakvog ograničenja na traku --- i pitanja odlučivosti.